\section{Deletion banks and necessary counterexample structure}
\label{sec:comb-banks}

Throughout this section, let $A\subseteq\{2,\ldots,M\}$ be a set of
$n$ distinct integers and let
\[
 I=\{x+1,\ldots,x+M\},\qquad x\in\mathbb Z_{\ge0}.
\]
The ambient length $M$ is retained when sources are deleted; it need
not equal the maximum of a smaller source set. Write
\[
 D_p(A)=\sum_{a\in A}v_p(a),\qquad
 N_I(d)=\#\{b\in I:d\mid b\},
\]
and let $\tau(A,I)$ be the minimum cardinality of a subset whose
product is divisible by $\prod_{a\in A}a$. We use $P^+(a)$ for the
largest prime factor, $\omega(a)$ for the number of distinct prime
factors, and $\Omega(a)$ for their number with multiplicity.

The following two facts will be used repeatedly:
\begin{equation}
 \left\lfloor\frac Md\right\rfloor\le N_I(d)
 \le\left\lceil\frac Md\right\rceil,
 \qquad
 \prod_{a\in A}a\mid M!\mid\prod_{b\in I}b.
 \label{eq:comb-supply}
\end{equation}
The divisibility $\prod A\mid M!$ uses distinctness of the sources;
$M!\mid\prod I$ follows from $\prod I=M!\binom{x+M}{M}$.
In particular, if a chosen subset has deficient $p$-valuation, some
unused interval position contains $p$. Adding that position reduces
the sum of the positive valuation deficits by at least one. We call
this a \emph{prime-unit repair}; different repairs always use
different positions.

\subsection{The local deletion hypothesis}

Let $K$ be a positive integer. For the necessary conditions below,
assume
\begin{equation}
 \tau(A,I)>K+1,
 \qquad
 \text{each }A\setminus\{a\}\text{ has a cover }B_a\subseteq I
 \text{ with }|B_a|\le K-1.
 \label{eq:comb-local}
\end{equation}
All conclusions apply to \emph{every} such cover $B_a$. For a
least-cardinality counterexample to the exact $2n$ bound, one takes
$K=2n-1$. A different choice needed for the asymptotic problem is
justified separately below.

\begin{lemma}[Forced deletion banks and a divisor gap]
\label{lem:comb-forcing}
Under \eqref{eq:comb-local}, for every $a\in A$ and every prime $q\mid a$,
\[
 C_q(a):=\{b\in I:a/q\mid b\}\subseteq B_a.
\]
Consequently every source prime is strictly smaller than $K$,
every source has $\Omega(a)\ge3$, and
\begin{equation}
 a>\frac MK,
 \qquad \frac aq>\frac MK\quad(q\mid a\text{ prime}).
 \label{eq:comb-quotient-size}
\end{equation}
No source has a divisor $d$ with
\begin{equation}
 K<d\le\frac MK.
 \label{eq:comb-divisor-gap}
\end{equation}
\end{lemma}

\begin{proof}
If $C_q(a)$ contained an unused position, add it to $B_a$. All
demands except possibly one $q$-unit are then supplied. One repair
finishes the cover in at most $K+1$ positions, a contradiction.
Since
$|C_q(a)|\ge\lfloor M/(a/q)\rfloor\ge q$, this inclusion gives
$q\le K-1$. It also gives $a/q>M/K$: otherwise the bank would have
at least $K$ members. An unused multiple of the whole source $a$
would finish the cover in one step, proving the analogous bound for
$a$. If $\Omega(a)\le2$, at most two prime-unit repairs would suffice,
which proves the remaining assertion about $\Omega$.

Finally, suppose $a=dv$ with $K<d\le M/K$. There are at least $K$
multiples of $d$, so one lies outside $B_a$. After choosing it, at
most $K$ positions are used. The number of multiples of $v$ is at
least $\lfloor M/v\rfloor\ge d\ge K+1$, so a second fresh position
supplies $v$. Their product supplies $a$, even if $d$ and $v$ are
not coprime, again contradicting \eqref{eq:comb-local}.
\end{proof}

The left endpoint in \eqref{eq:comb-divisor-gap} is strict. In
particular, this gap is empty when $M\le K^2$; it must not be used
as a nontrivial restriction in that range.

\subsection{Two-prime removal and a sufficient source class}

\begin{lemma}[A two-prime demand constraint]
\label{lem:comb-pair}
Assume \eqref{eq:comb-local}. Let $p,q$ be two prime occurrences in
$a\in A$, allowing $p=q$ only when $p^2\mid a$. If $pq\ge K$, then
\begin{equation}
 \begin{split}
 D_p(A)&>v_p\!\left(\prod_{b\in B_a}b\right)
 \ge v_p\!\left(\prod_{b\in C_q(a)}b\right)\\
 &\ge v_p(a/q)\left\lfloor\frac{Mq}{a}\right\rfloor
 \ge q\,v_p(a/q).
 \end{split}
 \label{eq:comb-pair-demand}
\end{equation}
\end{lemma}

\begin{proof}
Suppose instead that $B_a$ already supplies the full $p$-demand of
$A$. There are at least $pq\ge K$ multiples of $a/(pq)$ in $I$,
whereas $|B_a|\le K-1$. Choose one fresh multiple. If $p\ne q$,
every demand is now supplied except possibly one $q$-unit: the
$p$-demand was already complete, and the new position supplies the
other factors of $a$. One repair completes a cover in at most
$K+1$ positions. If $p=q$, the full $p$-demand was already complete
and even the last repair is unnecessary. This contradiction proves
the strict inequality. The others follow from the forced-bank
inclusion and \eqref{eq:comb-supply}.
\end{proof}

Bank inclusion alone does not imply a lower bound on the full
$p$-demand: selected positions could oversupply $p$ while serving
other primes. The fresh multiple of $a/(pq)$ is the additional
argument that makes \eqref{eq:comb-pair-demand} valid.

\begin{theorem}[A sufficient class at every interval scale]
\label{thm:comb-recognition}
Suppose that each $a\in A$ either has $\Omega(a)\le2$, or contains
two prime occurrences $p,q$ for which
\begin{equation}
 pq\ge2n-1,
 \qquad D_p(A)\le q\,v_p(a/q).
 \label{eq:comb-recognition}
\end{equation}
Then $\tau(A,I)\le2n$ for every positive interval of length
$M\ge\max A$. No bound on $\Omega(a)$ is required in the second case.
\end{theorem}

\begin{proof}
Induct on $n$, keeping the ambient length fixed. The case $n=1$
follows by choosing a multiple of its source. The hypothesis is
hereditary: deleting sources decreases both the demands and the
threshold $2n-1$. Obtain a cover $B$ of $A\setminus\{a\}$ using
at most $2n-2$ positions. If $\Omega(a)\le2$, use prime-unit repairs.

Otherwise choose witnesses $p,q$ in \eqref{eq:comb-recognition}.
If $C_q(a)$ has an unused position, add it and repair at most one
$q$-unit. If $C_q(a)\subseteq B$, then
$v_p(\prod B)\ge qv_p(a/q)\ge D_p(A)$. A fresh multiple of
$a/(pq)$ exists, and the construction in Lemma~\ref{lem:comb-pair}
again finishes in at most two additions.
\end{proof}

\subsection{All forced banks must fit simultaneously}

\begin{proposition}[Forced unions and their valuations]
\label{prop:comb-union}
Under \eqref{eq:comb-local}, put
\[
 U_a=\bigcup_{\substack{q\mid a\\q\ \mathrm{prime}}}C_q(a),
 \quad c(a)=1+\sum_{\substack{q\mid a\\q\ \mathrm{prime}}}(q-1),
 \quad t=\left\lfloor\frac Ma\right\rfloor.
\]
Then
\begin{equation}
 |U_a|=\sum_{\substack{q\mid a\\q\ \mathrm{prime}}}N_I(a/q)
       -(\omega(a)-1)N_I(a)\le K-1,
 \qquad t\,c(a)\le K-1.
 \label{eq:comb-union-size}
\end{equation}
For $p\mid a$, write $e_p=v_p(a)$ and define
\begin{equation}
 F_p(a)=t\bigl(e_pc(a)-(p-1)\bigr).
 \label{eq:comb-bank-valuation}
\end{equation}
Then $v_p(\prod U_a)\ge F_p(a)$, and necessarily
\begin{equation}
 \sum_{\substack{p\mid a\\p\ \mathrm{prime}}}
       \min\{e_p,(D_p(A)-F_p(a))_+\}\ge3.
 \label{eq:comb-residual-units}
\end{equation}
If $\Omega(a)=3$, the stronger equalities and inequalities
\begin{equation}
 v_p\!\left(\prod B_a\right)=D_p(A)-e_p,
 \qquad D_p(A)\ge F_p(a)+e_p\quad(p\mid a)
 \label{eq:comb-three-units}
\end{equation}
hold for every deletion cover $B_a$.
\end{proposition}

\begin{proof}
For distinct primes $p,q\mid a$, one has
$\operatorname{lcm}(a/p,a/q)=a$. Thus two different banks intersect
exactly in the multiples of $a$. A multiple of $a$ is counted
$\omega(a)$ times in the sum, and every other union member only
once. This proves the identity in \eqref{eq:comb-union-size}.

In a complete period of length $a$, there is the common residue
$0$, and for each prime $q\mid a$ there are $q-1$ additional
residues $ja/q$, $1\le j<q$. They are disjoint across the primes.
The period therefore contains $c(a)$ union members. Partitioning
the interval into $t$ complete periods and a remainder proves the
uniform cardinality bound.

For a fixed $p\mid a$, the $p-1$ nonzero residues from its own bank
have valuation exactly $e_p-1$. Every other union residue has
valuation at least $e_p$. This proves
\eqref{eq:comb-bank-valuation} on each complete period, and hence
the asserted lower bound on the whole interval.

The deletion cover supplies at least $D_p(A)-e_p$ units at every
prime of $a$, and it contains $U_a$. Its remaining $p$-deficit is
therefore at most $\min\{e_p,(D_p(A)-F_p(a))_+\}$. A sum at most
two would permit two repairs, proving \eqref{eq:comb-residual-units}.
When $\Omega(a)=3$, any surplus over $D_p(A)-e_p$ at even one prime
of $a$ would leave at most two units. Thus all those surpluses are
zero, which gives \eqref{eq:comb-three-units}.
\end{proof}

For example, for a squarefree triple $a=pqr$ the last bound gives
$D_p(A)\ge\lfloor M/a\rfloor(q+r-1)+1$, and its cyclic analogues.
These are necessary conditions, not a proof that such a deletion
cover exists.

\subsection{Saturation by a prime-power bank}

\begin{lemma}[Prime-power saturation]
\label{lem:comb-power-saturation}
Let $A$ consist of $n$ distinct sources in $[2,M]$. If $B\subseteq I$
contains every multiple of $p^e$, where $e\ge1$, and
$\lfloor M/p^e\rfloor\ge n$, then $v_p(\prod B)\ge D_p(A)$.
\end{lemma}

\begin{proof}
Counting source valuation levels gives
\[
 D_p(A)\le\sum_{j\ge1}\min\left\{n,
                         \left\lfloor\frac M{p^j}\right\rfloor\right\}.
\]
Every multiple of $p^e$ contributes its first $e$ units, and all
interval multiples of higher powers belong to the same bank. Hence
\[
 \begin{split}
 v_p\!\left(\prod B\right)
 &\ge eN_I(p^e)+\sum_{j>e}N_I(p^j)\\
 &\ge en+\sum_{j>e}\left\lfloor\frac M{p^j}\right\rfloor
 \ge D_p(A).
 \end{split}
\]
\end{proof}

\begin{corollary}
\label{cor:comb-power-exclusion}
Under \eqref{eq:comb-local}, no source has the form
$a=p^e q$, with $p,q$ prime, $e\ge1$, and $p^e\le M/n$.
The primes $p,q$ may coincide. In particular, such a source is
excluded whenever $q\ge n$.
\end{corollary}

\begin{proof}
Deleting the $q$ occurrence forces all multiples of $p^e$ into
$B_a$. Lemma~\ref{lem:comb-power-saturation} supplies the full
$p$-demand. Only one $q$-unit can remain; if $q=p$, nothing remains.
This contradicts \eqref{eq:comb-local}. Finally,
$q\ge n$ implies $p^e=a/q\le M/n$.
\end{proof}

The prime-power form is essential to this proof. It does not replace
the threshold $M/K$ by $M/n$ for an arbitrary composite cofactor.

\subsection{A small core and one remaining prime per source}

\begin{theorem}[Structure of a deletion obstruction]
\label{thm:comb-core}
Assume \eqref{eq:comb-local}. Among the source primes, define
\[
 E=\{p:p\le\sqrt K\}\ \cup\ \{p:D_p(A)>p\},
 \qquad R=\{\text{source primes}\}\setminus E.
\]
Every $r\in R$ is the unique largest prime of any source containing
it, and occurs there to exponent one. Thus such a source is $a=rd$,
where $d$ is $E$-smooth and $P^+(d)<r$, with $P^+(1)=1$.
Moreover
\begin{equation}
 |E|\le\lfloor\sqrt K\rfloor+\sqrt{2n\log_2M},
 \qquad R\subset(\sqrt K,K).
 \label{eq:comb-core-elementary}
\end{equation}
If $n\le K\le C_0n$ and $M<16n^2$, with fixed $C_0$, then as
$n\to\infty$,
\begin{equation}
 |E|=O_{C_0}\!\left(\sqrt{\frac n{\log n}}\right),\qquad
 \sum_{r\in R}\frac1r\le\log2+o(1),\qquad
 \sum_{\substack{p\in E\\p>\sqrt K}}\frac1p=o(1).
 \label{eq:comb-core-sharp}
\end{equation}
\end{theorem}

\begin{proof}
If a source containing $r\in R$ had another prime occurrence
$q\ge r$, then $rq>K$ and
$D_r(A)\le r\le qv_r(a/q)$, contradicting
Lemma~\ref{lem:comb-pair}. This proves both the unique-largest-prime
assertion and the fact that the remaining cofactor is $E$-smooth.
Lemma~\ref{lem:comb-forcing} gives the upper endpoint $r<K$.

Let $H=\{p:D_p(A)>p\}$ and $h=|H|$. Distinctness of its primes gives
\[
 \frac{h(h+1)}2\le\sum_{p\in H}p
 <\sum_pD_p(A)=\sum_{a\in A}\Omega(a)\le n\log_2M,
\]
which proves the elementary bound. For the sharper estimate use
the logarithmic budget
\begin{equation}
 \sum_{p\in H}p\log p
 <\sum_pD_p(A)\log p
 =\sum_{a\in A}\log a\le n\log M.
 \label{eq:comb-log-budget}
\end{equation}
Chebyshev's bound $\pi(t)\ll t/\log t$ implies that the $j$th prime
is at least a constant times $j\log(j+1)$. For completeness, the
usual bound follows by applying
$\prod_{m<p\le2m}p\mid\binom{2m}{m}<4^m$, summing over dyadic
intervals to bound $\sum_{p\le t}\log p$, and separating primes
below and above $\sqrt t$. It follows that the sum of $p\log p$
over any $h$ primes is at least
$c_1h^2\log^2(h+1)$ for an absolute $c_1>0$.
Together with \eqref{eq:comb-log-budget} and $\log M=O(\log n)$,
this gives $h=O(\sqrt{n/\log n})$: if $h<n^{1/4}$ it is immediate,
and otherwise $\log(h+1)\ge\frac14\log n$.
The remaining part of $E$ has at most
$\pi(\sqrt K)=O_{C_0}(\sqrt n/\log n)$ primes.

The prime reciprocal asymptotic
$\sum_{p\le t}1/p=\log\log t+B+o(1)$, for example from
Mertens' estimates \cite{TaoMertens}, gives
\[
 \sum_{\sqrt K<p<K}\frac1p=\log2+o(1).
\]
To control the large-prime part of $E$, put $T=\sqrt K\log n$.
The sum over all primes in $(\sqrt K,T]$ tends to zero because
$\log T/\log\sqrt K\to1$. For the rest, use
\eqref{eq:comb-log-budget}:
\[
 \sum_{\substack{p\in H\\p>T}}\frac1p
 \le\frac{n\log M}{T^2\log T}=O_{C_0}(\log^{-2}n).
\]
This proves all of \eqref{eq:comb-core-sharp}.
\end{proof}

The set $E$ is defined by source demands, not by the magnitudes of
dual weights. Its sublinear size does not by itself supply a
fractional cover or a rounding argument for its interaction with $R$.

\subsection{Why the reduction also applies to an asymptotic obstruction}

Let $g(n)$ denote the supremum of $\tau(A,I)$ over $n$ distinct
sources, with interval length $M=\max A$. A smallest exact-$2n$
counterexample does not automatically describe failure of the
asymptotic conjecture. The following argument provides the required
quantifiers.

\begin{proposition}[Linear budgets with an additive constant]
\label{prop:comb-asym-minimality}
Fix $2<c<10$. Suppose that for every $C\ge0$ there is an instance
with $\tau(A,I)>cn+C$. For each $C$, choose the smallest failing
cardinality $n_C$ and one such instance. As $C\to\infty$, these
instances satisfy the conclusions of Theorem~\ref{thm:comb-core}
with
\[
 K=\lfloor c(n_C-1)+C\rfloor+1,
 \qquad 2n_C-1\le K<10n_C,
 \qquad M<16n_C^2,
\]
and $n_C\to\infty$.
\end{proposition}

\begin{proof}
Minimality supplies a cover of every deletion with at most
$\lfloor c(n_C-1)+C\rfloor=K-1$ positions. If the deletion has
smaller maximum, apply its bound inside a contained interval of
that smaller length. Since $c\ge2$,
\[
 K+1\le\lfloor cn_C+C\rfloor<\tau(A,I).
\]
Thus \eqref{eq:comb-local} holds.

The known uniform bound $g(n)\le10n$ \cite{Wang} implies
$C<(10-c)n_C$, hence $n_C\to\infty$ and $K<10n_C$; the lower
bound on $K$ follows from $c\ge2$ and $C\ge0$.
The known long-interval theorem
$M\ge16n_C^2\Rightarrow\tau(A,I)\le2n_C$ \cite{Wang}
forces $M<16n_C^2$. Theorem~\ref{thm:comb-core} now applies.
\end{proof}

If the unrestricted asymptotic coefficient-two assertion fails,
one can choose a fixed $c>2$ below the offending limiting upper
ratio, and then no finite additive constant $C$ works. Thus the
proposition is a genuine necessary reduction for a possible
asymptotic obstruction family. It does not prove that the growing
core can be absorbed, nor identify it with the exceptional support
of an optimal dual vector.

\section{A complete reconstruction of the \texorpdfstring{$32n^{5/2}$}{32 n to the five-halves} argument}
\label{sec:comb-dot}

This section reconstructs the supplied argument in full, including
the demand step and the parameter restriction on which its
three-occurrence structure depends. Its range is already contained
in the stronger quadratic result of Wang \cite{Wang}; it is not
presented as a new best long-interval bound.

\begin{theorem}
\label{thm:comb-dot}
For every set $A$ of $n$ distinct integers in $[2,M]$ and every
positive interval $I$ of length $M$,
\[
 M\ge32n^{5/2}\quad\Longrightarrow\quad\tau(A,I)\le2n.
\]
\end{theorem}

\begin{proof}
Induct on $n$ at fixed ambient length $M$. The case $n=1$ is
immediate. The induction hypothesis is available for a deletion
because $M\ge32n^{5/2}\ge32(n-1)^{5/2}$. Suppose a counterexample
exists with $n\ge2$, and set
\[
 K=2n-1,\qquad c=\frac M{K^3},\qquad L=\frac MK=cK^2.
\]
The local deletion hypothesis holds. In particular, every source
prime is below $K$, every prime-deletion quotient exceeds $L$,
and no source divisor lies in $(K,L]$.

\paragraph{The parameter range.}
If $c\ge1$, every source exceeds $L\ge K^2$. Its least divisor
$d>K$ has all proper divisors at most $K$. Removing any prime
factor $p<K$ gives $d/p\le K$, so $d<K^2\le L$, a contradiction.
We may therefore assume $0<c<1$. The numerical hypothesis gives
\begin{equation}
 c\ge\frac{32n^{5/2}}{(2n-1)^3}>\frac4{\sqrt n},
 \qquad c^2K\ge1.
 \label{eq:comb-dot-range}
\end{equation}
All structural steps that follow use the explicit restrictions
$0<c<1$ and $c^2K\ge1$. They imply
$cK>1$ and $c^{3/2}K>1$.

\paragraph{Finding a large prime occurrence.}
Choose the least divisor $d>K$ of a source. If $\Omega(d)\ge3$,
its smallest prime factor is at most $d^{1/3}$; removing it leaves
a proper divisor at least $d^{2/3}$. Thus
\[
 d\le K^{3/2}\le cK^2=L,
\]
contradicting the divisor gap. The divisor is not prime because
all source primes are below $K$. Hence $d$ is a product of two
prime occurrences, and its larger prime $q$ satisfies
\begin{equation}
 q>\sqrt c\,K.
 \label{eq:comb-dot-large-prime}
\end{equation}

\paragraph{The small part.}
Put $U=1/\sqrt c$. There is no source divisor $t$ with
$U<t\le cK$. Indeed, the prime in
\eqref{eq:comb-dot-large-prime} is larger than $t$, since
$\sqrt cK>cK$. Therefore $qt$ divides the source and
\[
 K<qt<cK^2=L,
\]
again impossible. Every source prime is consequently either at
most $U$ or larger than $cK$.

Let $u$ be the full part of a source supported on primes at most
$U$. Then $u\le U$. Otherwise its least divisor $t>U$ is at most
$U^2=1/c$: remove a prime factor at most $U$ and use minimality.
But $1/c\le cK$ by \eqref{eq:comb-dot-range}, contradicting the
preceding forbidden interval.

\paragraph{Exactly three large occurrences.}
Call a prime occurrence larger than $cK$ large. Any two large
occurrences have product greater than $c^2K^2\ge K$. The divisor
gap therefore forces their product above $L$. Four such
occurrences would have product greater than
$L^2=M(cK)>M$, so a source has at most three.

It cannot have at most two. Delete a prime occurrence so that at
most one large occurrence remains; if none is large, delete any
prime. The quotient is then at most
\[
 uK\le UK=\frac K{\sqrt c}<cK^2=L,
\]
where $c^{3/2}K>1$ was used. This contradicts
\eqref{eq:comb-quotient-size}. Thus every source has the canonical
form
\begin{equation}
 a=u p q r,\qquad u\le U,
 \label{eq:comb-dot-triples}
\end{equation}
where $p,q,r$ are three large prime occurrences, with repetition
allowed. The small part $u$ need not have a bounded number of
prime factors.

\paragraph{Full demands at the large primes.}
Fix any large occurrence $p$ in a source
\eqref{eq:comb-dot-triples}. The other two have product above
$L=cK^2$, so one of them, say $q$, is larger than $\sqrt cK$.
Also $pq>L>K$. Lemma~\ref{lem:comb-pair} gives
\begin{equation}
 D_p(A)>qv_p(a/q)\ge q>\sqrt cK.
 \label{eq:comb-dot-demand}
\end{equation}
This remains valid when occurrences have equal prime values.
It is the two-prime removal argument, rather than bank inclusion
alone, that justifies \eqref{eq:comb-dot-demand}.

Let $h$ be the number of distinct large primes across the sources.
There are exactly $3n$ large-prime occurrences in total, so
\[
 h<\frac{3n}{\sqrt cK}\le\frac2{\sqrt c}=2U,
\]
where $3n\le2(2n-1)$ for $n\ge2$.

\paragraph{Counting the distinct sources.}
Unique factorization makes \eqref{eq:comb-dot-triples} an injection
into the pairs consisting of an integer $u\le U$ and a multiset
of three primes from the $h$ available large primes. Therefore
\[
 n\le\lfloor U\rfloor\binom{h+2}{3}
   <\frac{U(4U)^3}{6}
   =\frac{64}{6c^2}
   <\frac23n.
\]
Here $U>1$, $h<2U$, and $c>4/\sqrt n$ were used. The contradiction
completes the induction.
\end{proof}

The same three-occurrence argument cannot simply be continued to
$M=Cn^2$ with fixed $C$. Then $c=M/(2n-1)^3$ is of order $1/n$,
and the essential condition $c^2K\ge1$ fails for large $n$.
A quadratic proof therefore requires a different argument; the
independently checked theorem $M\ge16n^2\Rightarrow\tau\le2n$
\cite{Wang} supplies that range. None of the necessary conditions
above resolves the remaining unrestricted exact or asymptotic
problem.
